\documentclass[10pt,a4paper]{amsart}
\usepackage[margin=19mm]{geometry}
\usepackage{amsmath,amssymb,mathtools}
\usepackage[scr=rsfs,cal=euler]{mathalfa}
\usepackage{xfrac}
\usepackage[unicode,hidelinks,bookmarks=false]{hyperref}
\newcommand{\ie}{i.e.\ }
\newcommand{\cf}{cf.\ }
\newcommand{\resp}{respectively\ }
\newcommand{\Subsection}{\subsection}
\providecommand{\nobreakdash}{\nobreak\hbox{-}}
\setlength{\emergencystretch}{2em}
\multlinegap0pt
\usepackage{bm}
\usepackage{bbm}
\usepackage{dsfont}
\usepackage{xspace}
\usepackage{cancel}
\usepackage{mathtools}
\usepackage{cleveref}
\usepackage{amsthm}
\makeatletter
\@ifundefined{lemma}{\newtheorem{lemma}{Lemma}}{}
\makeatother
\newcommand{\RR}{\mathbb{R}}
\newcommand{\ZZ}{\mathbb{Z}}
\title[Sparse maps]{Sparse maps}
\author{Elia Portnoy}
\date{Source: arXiv:2407.05137v3 (CC BY 4.0)}
\begin{document}
\maketitle

\noindent\emph{Source and licence.} Sparse maps. Version arXiv:2407.05137v3 (CC BY 4.0), distributed under the Creative Commons Attribution 4.0 International licence, as recorded in the arXiv licence metadata for this exact version and shown on the abstract page https://arxiv.org/abs/2407.05137v3. Full rights notice: licenses/arxiv-2407.05137-CC-BY-4.0.txt.\par\smallskip

\noindent\textit{Editorial scope.} Counted unit: the Lemma whose source label is \texttt{extension} in the article (source0.tex lines 269--277, source label \texttt{extension}), delivered together with the article's own complete original proof of it (source0.tex lines 279--299, one \texttt{proof} environment of 21 lines). No mathematics is written, repaired or re-derived: statement and proof are the article's own text line for line. Required context retained uncounted: tight (lines 184--202). Every definition, display and label of those lines is retained unchanged. Omitted from this package and not redistributed: 1--183, 203--268, 278--278, 300--418. Nothing inside a retained excerpt is omitted or altered: every retained body is delivered exactly as the article's lines are written, so no omission receipt of any kind accompanies this package. Citations: the retained excerpts carry no citation command at all, so no bibliography entry of the article is transcribed and no reference is redistributed. Reference adaptation: none was needed and none was made; every reference inside the delivered unit resolves inside it, so no unresolved reference or citation is produced. Printed slips kept literal: no printed word, symbol, hypothesis or number of the counted statement or of its proof was altered. Layout and class: the article's own class is \texttt{article} with packages including amsmath, amsfonts, amssymb, amsthm, mathtools, verbatim, multicol, bbold, enumitem, graphicx, geometry, fancyhdr, cleveref, indentfirst; the wrapper is the lane's pinned \texttt{amsart} editorial wrapper at 10pt on A4, which restates the theorem-like environments the retained text needs and adds the article's own macro definitions that the retained text needs (\texttt{\textbackslash RR}, \texttt{\textbackslash ZZ}), with extra wrapper packages (bm, bbm, dsfont, xspace, cancel, mathtools, cleveref). The delivered numbering is the unit's own. Assets: no retained span contains a figure environment, a graphics-inclusion command, a PDF-page inclusion, a table or a TikZ picture; no image file is copied into this package, the package declares no asset dependency and no image file is redistributed with it. Licence: version arXiv:2407.05137v3 of this article is distributed under the Creative Commons Attribution 4.0 International licence, as recorded in the arXiv licence metadata for this exact version and shown on the abstract page https://arxiv.org/abs/2407.05137v3. A full rights notice is delivered at licenses/arxiv-2407.05137-CC-BY-4.0.txt. Characters: every retained excerpt is pure ASCII, so the delivered engine, XeLaTeX with its default Latin Modern text font, reports no missing character.
	\begin{lemma} \label{tight} Fix  dimensions $0 \le d \le m < n$, $d \le q \le n-1$, and an integer $S > 0$. Let $\{\sigma_1, \sigma_2 \ldots \sigma_N\}$ be a set of $d$-simplices and set $X = \bigsqcup_{i=1}^N \partial \sigma_i$. Suppose for some $R > 0$ we are given a map
		
		$$g: X \to [0,R]^{n-1} \times \{0\} \subset \RR^n$$
		
		\noindent which is $(m-1)$-sparse map with overlap $S$. Furthermore, suppose that $g$ is $q$-tight with respect to some $q$-plane $\Pi \subset \RR^{n-1} \times \{0\}$ and suppose that $g(\partial \sigma_i)$ does not lie in an $(d-1)$-plane for $1 \le i \le N$. Then there is a map
		
		$$G: X \times [0,2] \to  [0, R]^{n} \subset \RR^n$$
		
		\noindent so that the following conditions hold
		
		\begin{enumerate}
			\item $G$ extends $g$ in the sense that, $G|_{X \times \{0\}} = g$
			\item $G(X \times \{2\}) \subset [0,R]^{n-2} \times \{0\} \times [0,R]$
			\item $G|_{X \times \{2\}}$ is $(q-1)$-tight with respect to a $(q-1)$-plane contained in $\RR^{n-2} \times \{0\} \times \RR$.
			\item $G|_{X \times \{2\}}$ is $(m-1)$-sparse with overlap $S'(n,S)$ as a map into $\RR^{n-2} \times \{0\} \times \RR$. Here $S'(n,S)$ is a constant which only depends on $n$ and $S$.
			\item $G$ is $m$-sparse with overlap $S'(n,S)$.
		\end{enumerate} 
		
	\end{lemma}

	\begin{lemma} \label{extension} Fix three dimensions $1 \le d \le m < n$ and let $Y$ be a $d$-dimensional simplicial complex with at most $V$ vertices, so that each vertex lies in at most $D$ simplices. Suppose that for some $R>0$ we have a map,
		$$f: Y(d-1) \to [0, R]^{n-1} \times \{0\} \subset \RR^{n}$$
		
		\noindent which is $(m-1)$-sparse with overlap $S$. Then there is an extension of $f$,
		
		$$F: Y \to [0, R]^n  \subset \RR^{n}$$
		
		\noindent which is $m$-sparse with overlap $S'(n,D,S)$, where $S'(n,D,S)$ only depends on $n, D$ and $S$.
	\end{lemma}

	\begin{proof} Let $\sigma_1, \sigma_2 \ldots \sigma_N$ be the $d$-simplices of $Y$ and let $X = \bigsqcup_{i=1}^N \partial \sigma_i$. We will view $f$ as a map from, this is because each $\partial \sigma_i$ is a subset of both $X$ and of $Y$. Note that $f$ is $(m-1)$-sparse with overlap $DS$ when viewed as a map from $X$. Since the constant $S'(n,D,S)$ is allowed to depend on $S$ and $D$, it suffices to assume that $Y = \bigsqcup_{i=1}^N \sigma_i$ and $Y(d-1)$ is the $X$ just defined. This is because we have already fixed our map on the boundaries of these simplices. We will construct the map $F$ as a concatenation of $n-1$ maps. The  $i^{th}$ map will be $F_i: X \times [0,2] \to [0,R]^n$ and the maps will satisfy the following properties
		
		\begin{enumerate}
			\item $F_{n-1}$ extends $f$ and for each $i$, $F_i$ extends $F_{i-1}|_{X \times \{2\}}$. In other words, 
			$$F_i|_{X \times \{0\}} = F_{i-1}|_{X \times \{2\}}$$
			\item $F_i|_{X \times \{2\}}$ is $(i-1)$-tight and has image in some side of the cube $[0,R]^n$.
			\item $F_i$ is $m$-sparse with constant $S_i(n,D,S)$, which only depends on $n,D$ and $S$.
		\end{enumerate}
		
		We now explain how to define these maps. Define $F_{n-1}$ by using \cref{tight} to extend the given $(n-1)$-tight map $f$. The map $f$ is $(n-1)$-tight simply because its image is contained in a $(n-1)$-plane. \cref{tight} guarantees that $F_{n-1}|_{X \times \{2\}}$ is $(n-2)$-tight, and that $F_{n-1}$ is $m$-sparse with overlap $S_{n-1}(n,D,S)$ which only depends on $n,D$ and $S$. Next, we define $F_{n-2}$ by using \cref{tight} to extend the $(n-2)$-tight map $F_1|_{X \times \{2\}}$. Strictly speaking, to apply \cref{tight}, we need to rotate the coordinate system so that the image of $F_{n-1}|_{X \times \{2\}}$ lies inside $[0,R]^{n-1} \times \{0\}$. \cref{tight} guarantees that $F_{n-2}|_{X \times \{2\}}$ is $(n-3)$-tight, and that $F_{n-2}$ is $m$-sparse with overlap $S_{n-2}(n,D,S)$ which only depends on $n,D$ and $S_{n-1}(n,D,S)$. We can continue in this way to define each $F_i$. Specifically, to define $F_{i-1}$ apply \cref{tight} to the $(i-1)$-tight map $F_i|_{X \times \{2\}}$, after possibly rotating the coordinate system. By \cref{tight}, $F_{i-1}$ is $m$-sparse with overlap $S_{i-1}(n,D,S)$, which depends on $n,D$ an $S_i(n,D,S)$. But since $i < n$, we see that each constant $S_{i}(n,D,S)$ only depends on $n,D$ and $S$. Also notice that $F_{1}|_{X \times \{2\}}$ is $0$-tight. So for each $d$-simplex $\sigma$ in $Y$, $F_{n-1}(\partial \sigma \times \{2\})$ is a point in $\ZZ^n$. 
		
		We now concatenate the $F_i$'s to define $F$. This concatenation can be written out explicitly if we parametrize each $d$-simplex $\sigma$ in $Y$ as
		
		$$\sigma = \{(x, t): x \in \partial \sigma, t \in [0, n-1]\}/ \{(x, 0): x \in \partial \sigma\}$$
		
		\noindent For $1 \le i \le n-1$, $x \in \partial \sigma$ and $t \in [0,1]$ define
		
		$$F(x, i - t) = F_i(x, 2t)$$
		
		\noindent This map is well defined since $F(\partial \sigma \times \{0\}) = F_{1}(\partial \sigma \times \{2\})$ is a point in $\ZZ^n$. Since each $F_i$ was obtained using \cref{tight}, the image of $F$ lies in $[0,R]^n$ and $F$ is $m$-sparse with overlap $S'(n,D,S)$, which is some constant that depends only on $n, D$ and $S$.\footnote{We can be a bit more specific about the constant $S'(n,D,S)$. Let $\alpha$ be a tower of exponentials with base 2 and height $n$. Then $S'(n,D,S)$ is very roughly bounded by $(nDS)^{\alpha}$. This constant seems far from optimal, but we did not attempt to improve it in this paper.}
	\end{proof}

\end{document}
